% Copy this file to start a new note. Builds with pdflatex (the default recipe).
\documentclass[11pt]{article}
\usepackage{quicknotes}

\title{Data Science 6}
\author{Mark}
\date{}

\begin{document}
\maketitle

% \lecture{number}{date}{topic}
\lecture{6}{Monday, June 30, 2026}{Day 6}


\section{Apply}

Using the apply function we can reuse a function and apply it to all the columns: See example documentation. Does not change values of the table. ex: \textbf{weather.apply(celsisus, 'High')}. 
Returns an array can simply pass the function too. Order of params matters too. 
\begin{note}
    \itshape To make change use with\_column and pass the name and table.apply func.
\end{note}

\begin{note}
    \itshape What we could see in exam: \textbf{sat.with\_column('Avg. Score', sat.apply(avg\_score, 'Critical Reading', 'Math', 'Writing'))}
\end{note}

\section{Scope}

\begin{definition}
    \textbf{Scope} is essentially the area in a program where a name can be acessed.
\end{definition}

Recall global scope vs function/local scope. Always check this on mcq. Inside a local scope python always looks at local scope then global scope. 

\end{document}
